%% lists.tex -- Edge-case tests for the `list' parameter type, including
%% lists whose item type is itself a complex parameter.
\documentclass{article}
\usepackage{lua-placeholders}
\usepackage[autolanguage]{numprint}

\loadrecipe[\jobname]{lists-spec.yaml}
\loadpayload[\jobname]{lists-payload.yaml}

\setlength{\parindent}{0pt}

\begin{document}

    \section*{Case 1: list of strings (defaults from recipe)}

    \texttt{\textbackslash param} concatenates with \texttt{\textbackslash paramlistconjunction}:\\
    \param{strings}

    \medskip

    Iterating with \texttt{\textbackslash forlistitem} --- the csname takes one argument:

    \newcommand\itemBullet[1]{\textbullet\ #1\quad}
    \forlistitem{strings}{itemBullet}

    \section*{Case 2: list of numbers}

    \param{numbers}

    \medskip

    \forlistitem{numbers}{itemBullet}

    \section*{Case 3: list of objects --- type survives}

    Each item in \texttt{people} is an object with fields \texttt{name},
    \texttt{role}, \texttt{age}.  Because the item type isn't a primitive,
    \texttt{\textbackslash forlistitem} pushes each item's fields onto the
    context stack and the row macro accesses them by name --- with no
    arguments.

    \newcommand\peopleRow{\name{} (\role, \age) -- }
    \forlistitem{people}{peopleRow}

    \medskip

    The same iteration written so each person becomes a sentence using
    the bound field macros \texttt{\textbackslash name},
    \texttt{\textbackslash role}, \texttt{\textbackslash age}:

    \newcommand\peopleSentence{\name{} works as \role{} and is \age{} years old.\par}
    \forlistitem{people}{peopleSentence}

\end{document}
